\section{Especificações dos componentes da subestação}
\subsection{Para-raios}
Será utilizado o pára-raios modelo NLZP-1210, de óxido de Zinco polimérico
com desligador automático, conforme esquema da Norma da COSERN. Sendo conectados a
haste de aterramento da subestação através de cabo cobreado de bitola 50mm$^2$. Com as seguintes
especificações na Tabela 2.

% Table generated by Excel2LaTeX from sheet 'Plan1'
\begin{table}[htbp]
  \centering
    \begin{tabular}{cc}
    \toprule
    \textbf{Item } & \textbf{Especificação} \\
    \midrule
    Corrente de descarga nominal normalizada (8/20 $\mu$s) & 10 kA \\
    Tensão nominal (60 Hz, valor eficaz) & 12 kV \\
    Tensão disruptiva de impulso atmosférico normalizado (1,2/50 $\mu$s) & 44.4 kV  \\
    Tensão disruptiva à freqüência industrial & 18 kV \\
    Tensão disruptiva máxima de manobra (valor eficaz) & 32 kV \\
    \bottomrule
    \end{tabular}%
  \caption{Especificações para-raios NLZP-1510}
\end{table}%


\subsection{Chave fusível}
2 Elos-fusível 30k com chave fusível de corrente nominal 100 A, tensão nominal 13,8 kV, classe de
tensão 15 kV, NBI 95 kV e corrente de interrupção 4 kA, 1 Elo-fusível 65k com chave fusível de corrente nominal 100A, tensão nominal de 13,8kV, classe de tensão 15kv, NBI 95kV.

\subsection{Terminação primária}
Terminação primária unipolar, uso externo, do tipo composto termocontrátil, para cabo de 50 mm$^2$
com isolamento XLPE, tensão nominal de 15 kV, corrente nominal de 100 A, tensão suportável de impulso
95 kV,.

\subsection{Cabo de alta tensão}
O Cabo de cobre isolado, unipolar, classe de isolação de 12 /20 kV, em EPR e seção transversal de
50 mm$^2$.
\subsection{Cabo de baixa tensão}
Os cabos unipolares com isolação EPR 0,6/1,0 kV, marca Afumex, capa de proteção PVC e seção nominal 400 mm$^2$, método de instalação E para os transformadores de 750 kVA.
Serão usados para cada fase dois condutores de cobre com a seção especificada.
%http://br.prysmiangroup.com/br/files/dimensionamento_bt.pdf
\subsection{Chave seccionadora tripolar}
Três chaves fusíveis tripolares de abertura simultânea de uso interno e acionamento manual
através de alavanca de manobra, operação sem carga, corrente nominal 100 A, com corrente de curta
duração para efeito térmico de 10 kA, classe tensão 15 kV e para efeito dinâmico de 20 kA.
\subsection{Medição}
A caixa de medição será posicionada no interior da subestação com dispositivo para o lacre da concessionária conforme descrito pela Norma da COSERN.

A medição será feita a três elementos através de eletrodutos aparentes e de aço galvanizado.

Devido a capacidade instalada de 1500 kVA a medição será feita pela concessionária no circuito primário, por meio de 3 transformador de potencial de medição (TP) classe de isolamento 15 kV, grupo 
de ligação 2, tensão secundária 115 V, classe de exatidão de 0,3$\%$  e 3 transformadores de corrente (TC) classe de isolamento 15 kV, classe exatidão 0,3$\%$, fator térmico FT$=1,5$, uso interno e relação de transformação 50-5 A, especificações conforme a Norma da COSERN.

Os TP's e TC's serão instalados em suporte apropriado, conforme Norma da concessionária.

\subsection{Barramento de alta tensão}
Conforme especificação da concessionária, o barramento de alta tensão terá barras de cobre com
dimensões 3/4” x 3/16” e capacidade de suportar correntes de curto-circuito assimétricas de até 3 kA.
%\subsection{Barramento de baixa tensão} % %
Conforme memorial de cálculo, serão instaladas barras de cobre com seção de 499 mm$^2$ e que
suporte correntes de curto-circuito de até 36 kA.
%\subsection{Bucha de passagem} % %
Buchas de passagem para uso interno-interno com tensão nominal de 15 kV, corrente de 100 A,
NBI 95 kV.
\subsection{Disjuntor de alta tensão}
Disjuntor tripolar, a vácuo, comando automático de abertura, tensão nominal 13,8 kV, corrente
nominal 630 A, capacidade de interrupção 500 MVA. O disjuntor deverá ser equipado com relés
de ação indireta (secundários) de fase e neutro, 50/50N e 51/51N, com características de tempo inverso e
com dispositivos de ação instantânea.
\subsection{Disjuntor de baixa tensão}
Para os disjuntores dos QGFs alimentados pelos transformadores de 750 kVA serão utilizados disjuntores tripolares com as seguintes características:
\begin{itemize}
\item Tipo: 3VT5 H  - Siemens
\item Relé térmico:630-1600
\item Relé magnético: 5.000-10.000 A
\item Classe de temperatura da unidade magnética: 80 ms
\item Capacidade de ruptura: 55 kA/380 V
\item Ajuste do relé térmico:${I}_{a}=1250$ A
\end{itemize}

\subsection{Eletroduto entre os transformadores e QGF} %felipe
Serão construídas em alvenaria e possuirão dimensões dadas 300 mm x 300 mm com cobertura
feita através de grade de ferro.


\subsection{Transformadores}
Para atender às necessidades atuais da subestação, foram adotados dois transformadores de 750 kVA:

\begin{itemize}
\item Potência: 750 kVA
\item Refrigeração: ONAN - Óleo Natural, Ar Natural - imerso em óleo isolante mineral
%\item Atmosfera: Não é Agressiva
%\item Proteção: IP00
\item Classe de Tensão (kV): 15 kV
\item Tensão Primária: 13,8/13,2/12,6/12,0/11,4/10,8/10,2 kV
\item Tensão Secundária: 380/220 V
\item Primário: Triângulo (delta)
\item Secundário: Estrela com neutro acessível
\item Deslocamento Angular: 30$^{\circ}$
\item Freqüência nominal: 60 Hz
\item Corrente de excitação: 1,3$\%$
\item Impedância a 75$^{\circ}\mathrm{C}$: 5$\%$
\item Comprimento (C) : 1520 mm
\item Largura (L) : 1985 mm
\item Altura (A) : 1855 mm
\item Peso: 2075 kg
\item A óleo
\end{itemize}


\subsection{Aterramento} %felipe
Os dados referentes ao projeto da malha de aterramento se encontra no memorial de cálculo.